\chapter{डॉप्लर प्रभाव}\label{ch:g12:doppler-effect}

एंबुलेंस सायरन बजाती पास आती है और बग़ल से निकल जाती है --- और ठीक उसी क्षण सुर
साफ़ सुनाई देने लायक़ नीचे सरक जाता है, जबकि चालक ने कुछ भी नहीं छुआ। सायरन
ईमानदार है; बेईमान तरंग है: गति उसे आगे दबा देती है और पीछे खींच देती
है। यह अध्याय उस दबाव को नापता है, फिर उसी विचार पर सवार होकर सायरनों से
\omterm{met:g10:signals-and-waves:echo}{रडार} बंदूक़ तक, धमनी में बहते रक्त तक, और आगे उस ग्रह तक जाता है जिसे
तारे के प्रकाश की एक हलकी डगमगाहट ने पकड़वा दिया।

\section{यह प्रभाव}

\begin{definition}[डॉप्लर प्रभाव]\label{def:g12:doppler-effect:doppler}
जब स्रोत और ग्राही एक-दूसरे के सापेक्ष चलते हों तब किसी तरंग की मिलती
आवृत्ति का बदल जाना \emph{डॉप्लर प्रभाव}\index{डॉप्लर प्रभाव} कहलाता
है: पास आते समय मिलती आवृत्ति $f'$ छोड़ी गई आवृत्ति $f$ से
ऊँची रहती है, दूर जाते समय नीची, और $f$ के बराबर केवल तभी जब उनकी दूरी पल
भर के लिए बदलनी बंद कर दे।
\end{definition}

\begin{remark}[तारत्व, प्रबलता नहीं]\label{rem:g12:doppler-effect:pitch}
एंबुलेंस के गुज़रते समय दो चीज़ें बदलती हैं। \emph{\omterm{def:g12:sound-acoustics:pitch-loudness}{प्रबलता}} इसलिए
चढ़ती-उतरती है कि तरंग दूरी के साथ फैल जाती है (\cref{ch:g12:sound-acoustics}); और
\emph{\omterm{def:g12:sound-acoustics:pitch-loudness}{तारत्व}} पास आने या दूर जाने की चाल के कारण खिसकता है, और गुज़रने
के क्षण झटके से गिर जाता है। \omterm{def:g12:sound-acoustics:pitch-loudness}{प्रबलता} बताती है ``कितनी दूर'';
\omterm{def:g12:sound-acoustics:pitch-loudness}{तारत्व} बताता है ``कितनी तेज़''।
\end{remark}

\section{चलता स्रोत: सटते तरंगाग्र}

असली चित्र यह है: शिखर स्रोत से निकल जाने के बाद माध्यम की संपत्ति हो जाता है।
हर शिखर $v$ तरंग-चाल से (\cref{ch:g12:mechanical-waves}) एक वृत्त की तरह फैलता है, और उस
बिंदु पर केंद्रित रहता है जहाँ वह \emph{छोड़ा} गया था --- माध्यम को न पता है न
परवाह कि स्रोत आगे बढ़ चुका है।

\begin{proposition}[चलता स्रोत]\label{prop:g12:doppler-effect:moving-source}
कोई स्रोत ऐसे माध्यम में सरल रेखा में चाल $v_s < v$ से चलते हुए
\omterm{def:g10:signals-and-waves:period}{आवर्तकाल} $T$ (आवृत्ति $f = 1/T$) के साथ शिखर छोड़ता है जिसमें
तरंग चाल $v$ से चलती है। माध्यम में \omterm{def:g10:relative-motion:frame}{विराम} में बैठा प्रेक्षक पाता
है
\[
\lambda' = (v \mp v_s)\,T, \qquad f' = f\,\frac{v}{v \mp v_s},
\]
जहाँ ऊपर वाले चिह्न पास आते स्रोत के लिए हैं और नीचे वाले दूर जाते स्रोत के
लिए।
\end{proposition}

\begin{proof}
दो शिखरों के बीच पहला $vT$ आगे बढ़ता है जबकि स्रोत $v_s T$। इसलिए आगे की ओर
दूसरा शिखर पहले से $v_s T$ \emph{पास} छूटता है, $\lambda' = (v - v_s)T$; और पीछे की ओर वही
$v_s T$ जुड़ जाता है। फिर शिखर \omterm{def:g10:relative-motion:frame}{विराम} में बैठे प्रेक्षक के पास से माध्यम की चाल
$v$ से गुज़रते हैं, हर $\lambda'/v$ \omterm{def:g10:orders-of-magnitude:unit}{सेकंड} में एक, इसलिए $f' = v/\lambda' = f v/(v \mp v_s)$।
\end{proof}

\begin{omfigure}
\begin{tikzpicture}[font=\small]
  \foreach \i/\r in {0/3.0, 1/2.25, 2/1.5, 3/0.75} {
    \draw[omDef] ({0.3*\i},0) circle ({\r});
    \fill[gray] ({0.3*\i},0) circle (1.1pt);
  }
  \fill (1.2,0) circle (2pt);
  \draw[-Stealth, omThm, thick] (1.2,0) -- (1.75,0)
    node[below right=-2pt] {$\vect{v}_s$};
  \draw[Stealth-Stealth, thick] (2.55,0) -- (3.0,0);
  \node[above] at (2.78,0.05) {$\lambda'$};
  \draw[Stealth-Stealth, thick] (-3.0,0) -- (-1.95,0)
    node[midway, above] {$\lambda'$};
  \node[gray] at (0.15,0.55) {\scriptsize पिछली स्थितियाँ};
\end{tikzpicture}

{\small दाईं ओर चलते स्रोत के \omterm{def:g12:mechanical-waves:wavefront}{तरंगाग्र}: हर वृत्त उसी बिंदु पर केंद्रित
है जहाँ वह छूटा था (धूसर बिंदु), इसलिए शिखर आगे सट जाते हैं, $\lambda' = (v - v_s)T$, और पीछे
खिंच जाते हैं, $\lambda' = (v + v_s)T$।}
\end{omfigure}

\begin{example}[एंबुलेंस, संख्याओं में]\label{ex:g12:doppler-effect:siren}
$f = \qty{700}{Hz}$ का सायरन $v_s = \qty{25}{m/s}$ से पास आता है; हवा में $v = \qty{340}{m/s}$ लीजिए (\cref{ch:g12:sound-acoustics})। पास आते
समय: $f' = 700 \times 340/315 \approx \qty{756}{Hz}$; दूर जाते समय: $f' = 700 \times 340/365 \approx \qty{652}{Hz}$। गुज़रने पर \omterm{def:g12:sound-acoustics:pitch-loudness}{तारत्व} \qty{104}{Hz} गिर जाता
है --- अनुपात $\approx 1.16$, यानी लगभग ढाई अर्धस्वर, जिसे कोई भी कान चूक नहीं सकता।
\end{example}

\begin{remark}[गति किसे नहीं बदलती]\label{rem:g12:doppler-effect:speed-unchanged}
तरंग अब भी $v$ से ही चलती है: चाल केवल माध्यम तय करता है (\cref{ch:g12:mechanical-waves})।
स्रोत की गति \emph{\omterm{def:g12:mechanical-waves:wavelength}{तरंगदैर्घ्य}} दबाती है, चाल कभी नहीं। और इस दबाव की
एक सीमा है: $v_s \to v$ होते-होते आगे के शिखर एक-दूसरे पर चढ़ने लगते हैं और $f' \to \infty$
--- यही तरंगाग्रों का वह जाम है जिसे पराध्वनिक विमान झटके के रूप में
घसीटता है।
\end{remark}

\section{चलता प्रेक्षक, और धीमी गति का शॉर्टकट}

\begin{proposition}[चलता प्रेक्षक]\label{prop:g12:doppler-effect:moving-observer}
कोई प्रेक्षक माध्यम में \omterm{def:g10:relative-motion:frame}{विराम} में बैठे स्रोत की ओर (या उससे दूर) सीधी दिशा में
चाल $v_o$ से चलता है। मिलती आवृत्ति है
\[ f' = f\left(1 \pm \frac{v_o}{v}\right) \quad (+\ \text{पास आना},\ -\ \text{दूर जाना}). \]
\end{proposition}

\begin{proof}
तरंग का ढर्रा ख़ुद अछूता रहता है: $\lambda = vT$ अंतराल के शिखर $v$ से चलते
रहते हैं। उनकी ओर दौड़ता प्रेक्षक शिखरों से सापेक्ष चाल $v + v_o$ पर मिलता है,
हर $\lambda/(v + v_o)$ \omterm{def:g10:orders-of-magnitude:unit}{सेकंड} में एक: $f' = (v + v_o)/\lambda = f(1 + v_o/v)$। दूर भागते समय $v - v_o$।
\end{proof}

\begin{definition}[त्रिज्य वेग]\label{def:g12:doppler-effect:radial}
किसी प्रेक्षक के सापेक्ष स्रोत का \emph{त्रिज्य वेग}\index{त्रिज्य वेग}
$v_r$ उनके सापेक्ष वेग का वह घटक है जो उन्हें जोड़ती रेखा के अनुदिश हो ---
यानी वह दर जिससे दूरी सिकुड़ती या बढ़ती है। आवृत्ति केवल यही घटक
खिसकाता है; दृष्टि-रेखा के \emph{आड़े} की गति इस कोटि पर कोई खिसकाव नहीं
देती।
\end{definition}

\begin{proposition}[धीमी गति: सबके लिए एक ही सूत्र]\label{prop:g12:doppler-effect:slow}
जब त्रिज्य चाल छोटी हो, $v_r \ll v$, तब यह मायने ही नहीं रखता कि चल कौन रहा है:
\[
\frac{\Delta f}{f} = \frac{f' - f}{f} \approx +\frac{v_r}{v} \ \text{(पास आना)},
\qquad
\frac{\Delta f}{f} \approx -\frac{v_r}{v} \ \text{(दूर जाना)}.
\]
\end{proposition}

\begin{proof}
प्रेक्षक वाले सूत्र तो पहले से ही ठीक-ठीक $1 \pm v_o/v$ हैं। स्रोत के लिए $x = v_s/v$
लिखिए: सर्वसमिका $\frac{1}{1 - x} = 1 + x + \frac{x^2}{1 - x}$ दिखाती है कि $f'/f$ $1 + x$ से $x^2$ की कोटि के एक पद
जितना भिन्न है --- और $x = 0.1$ के लिए यह $10\%$ के खिसकाव पर $1\%$ का सुधार भर
है। पहली कोटि तक चारों सूत्र सिमटकर $1 \pm v_r/v$ बन जाते हैं।
\end{proof}

\begin{omfigure}
\begin{tikzpicture}
\begin{axis}[omaxis, width=8.8cm, height=4.8cm,
    xlabel={$t$ (s)}, ylabel={$f'$ (Hz)},
    xmin=-8, xmax=8, ymin=408, ymax=472,
    ytick={420,440,460}, xtick={-8,-4,0,4,8}]
  \addplot[dashed, omExo, domain=-8:8, samples=2] {460.3};
  \addplot[dashed, omExo, domain=-8:8, samples=2] {421.4};
  \addplot[dotted, omProp, thick, domain=-8:8, samples=2] {440};
  \addplot[omDef, thick, domain=-8:8, samples=120]
    {440*340/(340 + 15*(15*x)/sqrt((15*x)^2 + 64))};
\end{axis}
\end{tikzpicture}

{\small माइक्रोफ़ोन से \qty{8}{m} दूर से \qty{15}{m/s} पर गुज़रते \qty{440}{Hz} के सायरन की
सुनाई देती आवृत्ति: $fv/(v \mp v_s)$ पर पठार, और सच्ची $f$ (बिंदुदार) लगभग
ठीक निकटतम पहुँच पर पार होती है, जहाँ गति विशुद्ध \omterm{def:g12:mechanical-waves:transverse}{अनुप्रस्थ} है।}
\end{omfigure}

\begin{method}[गुज़रना पढ़ना]\label{met:g12:doppler-effect:pass}
गुज़रते स्रोत का अभिलेख दो पठार दिखाता है: $f_1$ (पास आना) और $f_2$ (दूर
जाना)। तब
\[
f = \frac{2 f_1 f_2}{f_1 + f_2}, \qquad
v_s = v\,\frac{f_1 - f_2}{f_1 + f_2}.
\]
सचमुच $f_1 + f_2 = \frac{2fv^2}{v^2 - v_s^2}$ और $f_1 - f_2 = \frac{2fv\,v_s}{v^2 - v_s^2}$: $v_s$ के लिए भाग दीजिए, और $f$ के लिए $2f_1 f_2/(f_1 + f_2)$
निकाल लीजिए। दो आवृत्ति-पाठ्यांकों से सच्चा \omterm{def:g12:sound-acoustics:pitch-loudness}{तारत्व} और चाल,
दोनों मिल जाते हैं --- न कोई घड़ी, न कोई पैमाना।
\end{method}

\section{चाल नापती प्रतिध्वनियाँ}

किसी चलते लक्ष्य पर तरंग टकराइए और \omterm{def:g12:doppler-effect:doppler}{डॉप्लर प्रभाव} दो बार चोट
करता है।

\begin{proposition}[परावर्तन खिसकाव को दुगुना कर देता है]\label{prop:g12:doppler-effect:double}
आवृत्ति $f$ की तरंग सीधे सामने से चाल $u \ll v$ पर पास आते
परावर्तक की ओर भेजी जाती है। \omterm{rem:g10:signals-and-waves:honest}{प्रतिध्वनि} इस आवृत्ति से लौटती है:
\[ f' = f\,\frac{v + u}{v - u}, \qquad \text{अतः} \qquad \Delta f \approx \frac{2u}{v}\,f . \]
\end{proposition}

\begin{proof}
लगातार दो खिसकाव। चलते \emph{प्रेक्षक} के रूप में परावर्तक $f(1 + u/v)$ पाता है;
और जो वह पाता है उसे फिर छोड़ते हुए वह अब चलता \emph{स्रोत} है, इसलिए
\omterm{rem:g10:signals-and-waves:honest}{प्रतिध्वनि} $u \ll v$ के लिए $f' = f(1 + u/v)/(1 - u/v) \approx f(1 + 2u/v)$ पर पहुँचती है।
\end{proof}

\begin{example}[रडार बंदूक़]\label{ex:g12:doppler-effect:radar}
यातायात का \omterm{met:g10:signals-and-waves:echo}{रडार} $f = \qty{24.15}{GHz}$ ($v = c = \qty{3.00e8}{m/s}$) पर सूक्ष्मतरंगें छोड़ता है और
\omterm{rem:g10:signals-and-waves:honest}{प्रतिध्वनि} को जाती हुई तरंग पर चढ़ा देता है; और विस्पंद
(\cref{ch:g12:sound-acoustics}) सीधे कार की चाल के प्रति \unit{m/s} पर $\Delta f = 2uf/c \approx \qty{161}{Hz}$ सुना देता है
--- यानी \qty{24}{GHz} के वाहक में से तराशी गई ऐसी श्रव्य आवृत्ति जिसे नापना
बच्चों का खेल है। गुणक $2$ छोड़ा नहीं जा सकता: उसे भूल जाइए और हर चालक को
आधी रियायत मिल जाएगी।
\end{example}

\begin{example}[डॉप्लर पराश्रव्य]\label{ex:g12:doppler-effect:ultrasound}
कोई चिकित्सा-यंत्र किसी धमनी में $f = \qty{5.0}{MHz}$ का \omterm{def:g10:signals-and-waves:ultrasound}{पराश्रव्य} भेजता है (ऊतक
में $v \approx \qty{1540}{m/s}$); और लाल रक्त कोशिकाएँ उसे परावर्तित कर देती हैं। $u = \qty{0.50}{m/s}$ पर बहता
रक्त $\Delta f = 2 \times 0.50 \times \num{5.0e6}/1540 \approx
\qty{3.2}{kHz}$ लौटाता है --- यानी एक \emph{श्रव्य} विस्पंद: हृदय-रोग विशेषज्ञ
सचमुच हर धड़कन पर रक्त को तेज़ होते सुनता है, और फिर $\Delta f$ से $u$ नाप
लेता है।
\end{example}

\section{तारों के प्रकाश में डॉप्लर}

प्रकाश एक तरंग है (\cref{ch:g12:light-as-wave}), और उसका \omterm{def:g12:sound-acoustics:timbre}{वर्णक्रम} उन तीखी रेखाओं
से धारीदार है जिनके \omterm{def:g12:mechanical-waves:wavelength}{तरंगदैर्घ्य} हर रासायनिक तत्व तय कर देता है
(\cref{ch:g10:light-spectra}) --- यानी हर तारे पर छपा हुआ, प्रयोगशाला में अंशांकित एक पैमाना।

\begin{definition}[लाल विस्थापन और नील विस्थापन]\label{def:g12:doppler-effect:redshift}
जब किसी स्रोत की \omterm{def:g10:light-spectra:emission}{वर्णक्रम रेखाएँ} प्रयोगशाला की तुलना में \emph{बड़े}
तरंगदैर्घ्यों पर दिखें, तो \omterm{def:g12:sound-acoustics:timbre}{वर्णक्रम} \emph{लाल-विस्थापित}\index{लाल विस्थापन}
है; और जब वे \emph{छोटे} पर दिखें, तो \emph{नील-विस्थापित}\index{नील विस्थापन}
--- यानी \omterm{def:g10:light-spectra:wavelength}{दृश्य वर्णक्रम} के लाल या नीले सिरे की ओर खिसका हुआ।
\end{definition}

\begin{proposition}[प्रकाश का डॉप्लर खिसकाव]\label{prop:g12:doppler-effect:light}
त्रिज्य चाल $v_r \ll c$ से दूर जाते स्रोत के लिए हर \omterm{def:g12:mechanical-waves:wavelength}{तरंगदैर्घ्य} इतना खिंच
जाता है:
\[ \frac{\Delta \lambda}{\lambda} = \frac{\lambda' - \lambda}{\lambda} \approx \frac{v_r}{c} \quad \text{(लाल विस्थापन)}, \]
और पास आते स्रोत के लिए उतना ही सिकुड़ जाता है, $\Delta\lambda/\lambda \approx -v_r/c$ (\omterm{def:g12:doppler-effect:redshift}{नील
विस्थापन})।
\end{proposition}
\admitted

\begin{remark}[ईमानदार सूत्र कहाँ रहता है]\label{rem:g12:doppler-effect:honest}
प्रकाश को किसी माध्यम की ज़रूरत नहीं, इसलिए \omterm{def:g10:signals-and-waves:sound}{ध्वनि} वाली व्युत्पत्ति
यहाँ उतारी नहीं जा सकती --- ``माध्यम में \omterm{def:g10:relative-motion:frame}{विराम} में'' जैसी कोई चीज़ ही नहीं है।
ठीक सूत्र विशिष्ट आपेक्षिकता से आता है, इसी वर्ष आगे (\cref{ch:g12:special-relativity}); और सुधार
$(v_r/c)^2$ की कोटि के हैं, जो विमानों, तारों और पास की आकाशगंगाओं के लिए नगण्य
हैं। नीचे हम साफ़ अंतःकरण से धीमी गति वाला सूत्र काम में लेते हैं।
\end{remark}

\begin{omfigure}
\begin{tikzpicture}[font=\small]
  \foreach \y/\xl/\c/\t in {2.4/1.8/black/laboratory,
      1.2/0.6/omDef/approaching star, 0/3.72/omThm/receding star} {
    \fill[gray!18] (0,\y) rectangle (6,\y+0.55);
    \draw[\c, very thick] (\xl,\y) -- (\xl,\y+0.55);
    \node[right] at (6.15,\y+0.28) {\t};
  }
  \draw[dashed, gray] (1.8,-0.15) -- (1.8,2.95);
  \draw[-Stealth] (-0.2,-0.55) -- (6.4,-0.55)
    node[below left] {$\lambda$ (nm)};
  \foreach \x/\l in {0/656.0, 3/656.5, 6/657.0}
    \draw (\x,-0.48) -- (\x,-0.62) node[below] {\l};
\end{tikzpicture}

{\small तीन वर्णक्रमों में वही हाइड्रोजन रेखा: प्रयोगशाला में \qty{656.30}{nm} पर, पास
आते तारे के लिए \omterm{def:g12:doppler-effect:redshift}{नील-विस्थापित}, और दूर जाते तारे के लिए
\omterm{def:g12:doppler-effect:redshift}{लाल-विस्थापित}। बिंदुदार संदर्भ के सामने पढ़ा गया यह खिसकाव
\omterm{def:g12:doppler-effect:radial}{त्रिज्य वेग} दे देता है।}
\end{omfigure}

\begin{example}[तारे का भागना तौलना]\label{ex:g12:doppler-effect:vega}
चमकीले तारे वेगा के \omterm{def:g12:sound-acoustics:timbre}{वर्णक्रम} में प्रयोगशाला के \omterm{def:g12:mechanical-waves:wavelength}{तरंगदैर्घ्य}
\qty{656.30}{nm} वाली हाइड्रोजन रेखा \qty{656.27}{nm} पर नापी जाती है: $\Delta\lambda = \qty{-0.03}{nm}$, यानी
\omterm{def:g12:doppler-effect:redshift}{नील विस्थापन}, इसलिए वेगा $v_r = c\,|\Delta\lambda|/\lambda \approx \qty{1.4e4}{m/s}$ से पास आ रहा है --- यानी \qty{14}{km/s},
और वह भी बीस हज़ार में एक भाग के खिसकाव से पढ़ा हुआ।
\end{example}

\begin{example}[डगमगाते तारे: युग्म तारे और बहिर्ग्रह]\label{ex:g12:doppler-effect:wobble}
जब दो तारे एक-दूसरे की परिक्रमा करते हैं तो उनकी रेखाएँ
\omterm{def:g10:signals-and-waves:period}{आवर्ती} रूप से नीले और लाल की ओर झूलती रहती हैं --- यही
\emph{वर्णक्रमी युग्म तारा}\index{वर्णक्रमी युग्म तारा} है, जिसकी
कक्षाएँ जोड़े को अलग-अलग देखे बिना ही नाप ली जाती हैं। कोई ग्रह अपने
तारे के साथ यही करता है, बस बहुत मंद ढंग से: दोनों अपने साझा
\omterm{def:g11:forces-and-motion:com}{द्रव्यमान केंद्र} की परिक्रमा करते हैं, इसलिए तारे का \omterm{def:g12:doppler-effect:radial}{त्रिज्य
वेग} कुछ दर्जन \unit{m/s} जितना दोलन करता है --- यानी $\Delta\lambda/\lambda \sim 10^{-7}$ का खिसकाव। इसी को
पकड़कर सूर्य जैसे किसी तारे के चारों ओर पहला ग्रह खोजा गया था; और सप्ताहांत की
समस्या उसी खोज को, संख्याओं समेत, फिर से चला देती है।
\end{example}

\begin{remark}[आकाशगंगाओं का लाल विस्थापन]\label{rem:g12:doppler-effect:expansion}
दूर की आकाशगंगाओं के \omterm{def:g12:sound-acoustics:timbre}{वर्णक्रम} सब के सब \emph{लाल}-विस्थापित हैं, और जितनी दूर,
उतने ही अधिक: ब्रह्मांड की दूरियाँ खिंच रही हैं। पास की आकाशगंगाओं के लिए
$v_r = c\,\Delta\lambda/\lambda$ उनके दूर जाने की चाल पढ़ लेता है; और सबसे दूर वालों के लिए खिसकाव उससे
आगे निकल जाते हैं जितना धीमी गति का सूत्र ईमानदारी से सँभाल सकता है, इसलिए वह
हिसाब \cref{ch:g12:special-relativity} से आगे और विश्वविद्यालय के खंडों में फिर उठाया जाता है।
\end{remark}

\section{अभ्यास}

\begin{exercise}[$\star$]\label{exo:g12:doppler-effect:1}
दमकल का सायरन \qty{700}{Hz} पर बजता है और गाड़ी \qty{25}{m/s} ($v = \qty{340}{m/s}$) से चलती है।
आवृत्ति निकालिए जो सुनाई देगी (क) उसके आगे खड़े राहगीर को; (ख) पीछे
खड़े को। (ग) चालक को क्या सुनाई देता है?
\end{exercise}

\begin{exercise}[$\star$]\label{exo:g12:doppler-effect:2}
कोई मित्र कहता है: ``एंबुलेंस के जाते समय \omterm{def:g12:sound-acoustics:pitch-loudness}{तारत्व} इसलिए गिरता है कि
दूरी के साथ \omterm{def:g10:signals-and-waves:sound}{ध्वनि} कमज़ोर पड़ जाती है।'' गड्डमड्ड किए जा रहे दोनों
प्रभाव अलग कीजिए, और बताइए कि हर एक किस पर निर्भर करता है।
\end{exercise}

\begin{exercise}[$\star$]\label{exo:g12:doppler-effect:3}
कोई साइकिल-सवार \qty{700}{Hz} पर बजते ठहरे हुए सायरन की ओर सीधे \qty{6.0}{m/s} से चलता है।
उसे कौन-सी आवृत्ति सुनाई देगी? और सीधे दूर जाते समय? कौन-सा सूत्र
लागू होता है, और चलते-स्रोत वाला क्यों नहीं?
\end{exercise}

\begin{exercise}[$\star$]\label{exo:g12:doppler-effect:4}
रेलगाड़ी का भोंपू \qty{400}{Hz} पर बजता है जबकि गाड़ी \qty{30}{m/s} से चलती है। गाड़ी के
आगे, पीछे, और \omterm{def:g10:relative-motion:frame}{विराम} में \omterm{def:g12:mechanical-waves:wavelength}{तरंगदैर्घ्य} निकालिए। कौन-सा प्रेक्षक कौन-सा
पाता है?
\end{exercise}

\begin{exercise}[$\star$]\label{exo:g12:doppler-effect:5}
\qty{440}{Hz} का स्वर बजाता फ़ोन एक दौड़ने वाला \qty{5.0}{m/s} से लिए जा रहा है। धीमी गति
वाले सूत्र से उस खिसकाव का अनुमान लगाइए जो उस व्यक्ति को सुनाई देगा जिसकी ओर
दौड़ने वाला बढ़ रहा है। उसकी तुलना एक अर्धस्वर से कीजिए (आवृत्ति में
लगभग $6\%$): क्या कोई संगीतकार उसे पकड़ लेगा?
\end{exercise}

\begin{exercise}[$\star\star$]\label{exo:g12:doppler-effect:6}
कोई चमगादड़ \qty{50.0}{kHz} पर आवाज़ छोड़ते हुए \qty{6.0}{m/s} से सीधे दीवार की ओर उड़ता है।
दिखाइए कि उसे सुनाई देती \omterm{rem:g10:signals-and-waves:honest}{प्रतिध्वनि} $f' = f(v + v_b)/(v - v_b)$ पर लौटती है, फिर $f'$
और खिसकाव निकालिए। यह भेस बदले रडार-बंदूक़ का सूत्र ही क्यों है?
\end{exercise}

\begin{exercise}[$\star\star$]\label{exo:g12:doppler-effect:7}
\qty{24.15}{GHz} पर चलती \omterm{met:g10:signals-and-waves:echo}{रडार} बंदूक़ पास आती कार से \qty{4.0}{kHz} का विस्पंद नापती
है। कार की चाल \unit{m/s} और \unit{km/h} में निकालिए। ठीक \qty{50}{km/h} पर चलती कार कितना
विस्पंद देती?
\end{exercise}

\begin{exercise}[$\star\star$]\label{exo:g12:doppler-effect:8}
किसी धमनी के अनुदिश ताका गया \qty{5.0}{MHz} का डॉप्लर \omterm{def:g10:signals-and-waves:ultrasound}{पराश्रव्य} यंत्र (ऊतक
में $v = \qty{1540}{m/s}$) \qty{2.6}{kHz} का विस्पंद सुनता है। रक्त की चाल निकालिए। यंत्र को बहाव
के लंबवत रखने के बजाय उसके \emph{अनुदिश} तिरछा क्यों रखना पड़ता है?
\end{exercise}

\begin{exercise}[$\star\star$]\label{exo:g12:doppler-effect:9}
$f = \qty{1000}{Hz}$, $v = \qty{340}{m/s}$ और \qty{34}{m/s} की पास आने की चाल लीजिए। $f'$ ठीक-ठीक
निकालिए, जब (क) स्रोत चले, (ख) प्रेक्षक चले। दोनों की तुलना धीमी गति के
पूर्वानुमान $f(1 + v_r/v)$ से कीजिए और असहमति के आकार का कारण बताइए।
\end{exercise}

\begin{exercise}[$\star\star$]\label{exo:g12:doppler-effect:10}
माइक्रोफ़ोन के पास से गुज़रते स्कूटर के अभिलेख में \qty{465}{Hz} और \qty{415}{Hz} पर पठार
दिखते हैं। \cref{met:g12:doppler-effect:pass} का उपयोग करके भोंपू की सच्ची आवृत्ति और स्कूटर की
चाल \unit{km/h} में निकालिए। सच्ची आवृत्ति उनका औसत \emph{क्यों नहीं}
है?
\end{exercise}

\begin{exercise}[$\star\star$]\label{exo:g12:doppler-effect:11}
किसी तारे के \omterm{def:g12:sound-acoustics:timbre}{वर्णक्रम} में प्रयोगशाला के \omterm{def:g12:mechanical-waves:wavelength}{तरंगदैर्घ्य} \qty{656.3}{nm}
वाली हाइड्रोजन रेखा \qty{656.5}{nm} पर दिखती है। \omterm{def:g12:doppler-effect:redshift}{लाल विस्थापन} या
\omterm{def:g12:doppler-effect:redshift}{नील विस्थापन}? तारे का \omterm{def:g12:doppler-effect:radial}{त्रिज्य वेग} और दिशा निकालिए।
\end{exercise}

\begin{exercise}[$\star\star\star$]\label{exo:g12:doppler-effect:12}
कोई तारा किसी अनदेखे ग्रह के कारण \qty{55}{m/s} से डगमगाता है। \qty{500}{nm} की रेखा के
\omterm{def:g12:mechanical-waves:wavelength}{तरंगदैर्घ्य} के झूले का आयाम निकालिए, लंबाई के रूप में भी और
$\lambda$ के अंश के रूप में भी। \omterm{def:g10:orders-of-magnitude:oom}{परिमाण की कोटि}: ग्रह-खोज को $10^{7}$
में एक भाग तक स्थिर वर्णक्रममापियों की प्रतीक्षा क्यों करनी पड़ी?
\end{exercise}

\begin{exercise}[$\star\star\star$]\label{exo:g12:doppler-effect:13}
किसी आकाशगंगा के \omterm{def:g12:sound-acoustics:timbre}{वर्णक्रम} की हर रेखा $2.4\%$ जितनी खिंची मिलती है।
उसके दूर जाने की चाल निकालिए। खगोलविद पाते हैं कि आकाश की हर दिशा में ऐसी
चालें दूरी के समानुपात में बढ़ती हैं: इससे कौन-सा चित्र बनता है? क्या यहाँ धीमी
गति का सूत्र अब भी भरोसे लायक़ है?
\end{exercise}

\begin{exercise}[$\star\star\star$]\label{exo:g12:doppler-effect:14}
सूर्य घूमता है: उसकी चकती का एक किनारा हमारी ओर लगभग \qty{2.0}{km/s} से आता है जबकि
दूसरा उतनी ही तेज़ी से दूर जाता है। पूरी चकती से आते प्रकाश की \qty{656.3}{nm} रेखा का
यह क्या करता है --- उसे खिसका देता है या चौड़ा कर देता है? असर \unit{nm} में
निकालिए।
\end{exercise}

\begin{exercise}[$\star\star\star$]\label{exo:g12:doppler-effect:15}
\qty{40}{m/s} से चलती पुलिस की गाड़ी, जिसका सायरन \qty{700}{Hz} पर है, उसी दिशा में \qty{30}{m/s}
से चलते ट्रक का पीछा करती है। (क) ट्रक-चालक के लिए $f' = f(v - v_o)/(v - v_s)$ का औचित्य बताइए।
(ख) $f'$ निकालिए। (ग) दिखाइए कि ट्रक गाड़ी की चाल पकड़ ले तो खिसकाव लुप्त हो
जाएगा, और बताइए कि ऐसा होना ही क्यों चाहिए।
\end{exercise}

\section{समस्या: बहिर्ग्रह की खोज}

\begin{problem}[{सप्ताहांत समस्या --- बहिर्ग्रह की खोज: रेलगाड़ी का भोंपू कान
अंशांकित करता है, रडार बंदूक़ दुहरे खिसकाव का पूर्वाभ्यास कराती है, और रविवार
की रात तक पचास लाख में एक भाग जितना डगमगाता एक तारा उस ग्रह को तौल देता है
जिसे किसी ने कभी नहीं देखा}]
\label{pb:g12:doppler-effect:1}
खगोल-मंडली एक सप्ताहांत एक ही विचार पर बिताती है --- डॉप्लर खिसकाव --- पहले
पृथ्वी पर उसका पूर्वाभ्यास, फिर सूर्य जैसे किसी तारे पर उसका निशाना। आँकड़े:
हवा में \omterm{def:g10:signals-and-waves:sound}{ध्वनि} के लिए $v = \qty{340}{m/s}$, $c = \qty{3.00e8}{m/s}$, $G = \qty{6.67e-11}{N.m^2/kg^2}$; तारे का द्रव्यमान
$M = \qty{2.0e30}{kg}$। वृत्तीय कक्षाओं की यांत्रिकी से मान लीजिए (\cref{ch:g12:satellites-kepler}): त्रिज्या
$r$ की वृत्तीय कक्षा पर द्रव्यमान $m$ के ग्रह का
\omterm{def:g10:signals-and-waves:period}{आवर्तकाल} और चाल $r^3 = \frac{G M P^2}{4\pi^2}$ तथा $V = \frac{2\pi r}{P}$ से मिलते हैं, और साझा
\omterm{def:g11:forces-and-motion:com}{द्रव्यमान केंद्र} की परिक्रमा करता तारा चाल $K = \frac{m}{M}\,V$ से डगमगाता है।

\medskip
\noindent\textbf{भाग I --- शुक्रवार: फाटक वाली क्रॉसिंग।}
कोई रेलगाड़ी भोंपू बजाती अचर चाल से गुज़रती है; एक फ़ोन पास आते समय $f_1 = \qty{471}{Hz}$ और
गुज़र जाने के बाद $f_2 = \qty{411}{Hz}$ दर्ज करता है।
\begin{enumerate}
  \item सटते तरंगाग्रों के चित्र से समझाइए कि भोंपू के कभी न बदलने
        पर भी $f_1 > f_2$ क्यों होता है।
  \item $f_1$, $f_2$ को सच्ची आवृत्ति $f$ और गाड़ी की चाल
        $u$ से जोड़ते दोनों समीकरण लिखिए।
  \item दिखाइए कि $u = v\,\dfrac{f_1 - f_2}{f_1 + f_2}$, और उसे \unit{m/s} तथा \unit{km/h} में निकालिए।
  \item दिखाइए कि $f = \dfrac{2 f_1 f_2}{f_1 + f_2}$, और उसे निकालिए। $f$ $f_1$ और $f_2$ के औसत
        से थोड़ा \emph{नीचे} क्यों है?
  \item गाड़ी के आगे और पीछे के \omterm{def:g12:mechanical-waves:wavelength}{तरंगदैर्घ्य} और \omterm{def:g10:relative-motion:frame}{विराम} वाला
        \omterm{def:g12:mechanical-waves:wavelength}{तरंगदैर्घ्य} निकालिए।
\end{enumerate}

\medskip
\noindent\textbf{भाग II --- शनिवार: \omterm{met:g10:signals-and-waves:echo}{रडार} बंदूक़।}
यातायात दल अपना \omterm{met:g10:signals-and-waves:echo}{रडार} दिखाता है: $f = \qty{24.125}{GHz}$, और बंदूक़ जाती हुई
तरंग तथा \omterm{rem:g10:signals-and-waves:honest}{प्रतिध्वनि} के बीच का विस्पंद पढ़ती है।
\begin{enumerate}[resume]
  \item समझाइए कि चलती कार की \omterm{rem:g10:signals-and-waves:honest}{प्रतिध्वनि} \emph{दो बार} क्यों खिसकती
        है --- हर खिसकाव में कार की भूमिका बताइए।
  \item दिखाइए कि चाल $u \ll c$ पर चलती कार के लिए विस्पंद $\Delta f \approx \dfrac{2u}{c}\,f$ होता है।
  \item कोई कार $\Delta f = \qty{3.55}{kHz}$ लौटाती है: $u$ निकालिए।
  \item उसे \unit{km/h} में बदलिए और वहाँ लगी \qty{80}{km/h} की सीमा से तुलना कीजिए।
  \item \qty{1}{km/h} के अनुरूप कितने \omterm{def:g10:signals-and-waves:frequency}{हर्ट्ज़} विस्पंद बनते हैं? टिप्पणी
        कीजिए: \qty{24}{GHz} की तरंग के इतने छोटे सापेक्ष खिसकाव को
        नापना आसान किस बात से हो जाता है?
\end{enumerate}

\medskip
\noindent\textbf{भाग III --- रविवार: डगमगाता तारा।}
मंडली सूर्य जैसे किसी तारे का, रात-दर-रात नापा गया \omterm{def:g12:doppler-effect:radial}{त्रिज्य वेग} उतार
लेती है, जो नीचे आलेखित है।

\begin{omfigure}
\begin{tikzpicture}
\begin{axis}[omaxis, width=8.8cm, height=4.6cm,
    xlabel={$t$ (दिन)}, ylabel={$v_r$ (m/s)},
    xmin=0, xmax=8.4, ymin=-72, ymax=72,
    xtick={0,2.1,4.2,6.3,8.4},
    xticklabels={$0$,$2.1$,$4.2$,$6.3$,$8.4$}, ytick={-55,0,55}]
  \addplot[dashed, omExo, domain=0:8.4, samples=2] {55};
  \addplot[dashed, omExo, domain=0:8.4, samples=2] {-55};
  \addplot[omDef, thick, domain=0:8.4, samples=140] {55*sin(360*x/4.2)};
\end{axis}
\end{tikzpicture}

{\small बारह रातों में तारे का \omterm{def:g12:doppler-effect:radial}{त्रिज्य वेग}: एक साफ़ ज्यावक्र --- यानी
किनारे से देखी गई वृत्तीय कक्षा पर घूमते किसी पिंड की छाप।}
\end{omfigure}

\begin{enumerate}[resume]
  \item रात-दर-रात तारे के \emph{\omterm{def:g12:sound-acoustics:timbre}{वर्णक्रम}} में असल में नापा क्या
        जाता है? उसे $v_r$ में बदलने वाला संबंध लिखिए।
  \item वक्र से डगमगाहट का आयाम $K$ और \omterm{def:g10:signals-and-waves:period}{आवर्तकाल} $P$
        पढ़िए।
  \item \qty{550}{nm} की रेखा का उससे बनता झूला $\Delta\lambda$ और $\Delta\lambda/\lambda$ निकालिए। भाग I
        से तुलना कीजिए: यह माप रेलगाड़ी वाली माप से कितनी गुना कठिन है?
  \item कोई अनदेखा ग्रह तारे को डगमगाता ही क्यों है? ज्यावक्र जैसा आकार
        कक्षा के बारे में क्या बताता है?
  \item यह विधि तारे की गति का केवल \emph{त्रिज्य} हिस्सा ही क्यों नापती है,
        और यदि कक्षा हमारे सामने चपटी पड़ी होती तो क्या दिखता?
\end{enumerate}

\medskip
\noindent\textbf{भाग IV --- रविवार की रात: अदृश्य को तौलना।}
कक्षा को वृत्तीय और किनारे से देखी हुई मानिए, ताकि $K$ तारे की
पूरी कक्षीय चाल हो।
\begin{enumerate}[resume]
  \item $P$ और मानी हुई संगति $r^3 = G M P^2/(4\pi^2)$ से ग्रह की कक्षीय त्रिज्या $r$
        निकालिए। उसकी तुलना पृथ्वी--सूर्य दूरी \qty{1.5e11}{m} से कीजिए।
  \item ग्रह की कक्षीय चाल $V = 2\pi r/P$ निकालिए।
  \item $K = (m/M)\,V$ से ग्रह का द्रव्यमान $m$ निकालिए।
  \item $m$ की तुलना बृहस्पति ($\qty{1.9e27}{kg}$) और पृथ्वी ($\qty{6.0e24}{kg}$) से कीजिए, और
        $r$ की तुलना सूर्य--बुध दूरी ($\qty{5.8e10}{m}$) से। यह कैसा संसार है?
  \item यदि कक्षा किनारे से नहीं, तिरछी होती तो सच्चा द्रव्यमान आपके
        मान से बड़ा निकलता या छोटा? एक वाक्य में निष्कर्ष लिखिए: पचास लाख में
        एक भाग के \omterm{def:g10:signals-and-waves:period}{आवर्ती} खिसकाव ने अभी-अभी क्या सौंप दिया?
\end{enumerate}
\end{problem}
